%bảng với độ rộng cột xác định
\usepackage{array}
\newcommand{\PreserveBackslash}[1]{\let\temp=\\#1\let\\=\temp}
\newcolumntype{C}[1]{>{\PreserveBackslash\centering}p{#1}}
\newcolumntype{R}[1]{>{\PreserveBackslash\raggedleft}p{#1}}
\newcolumntype{L}[1]{>{\PreserveBackslash\raggedright}p{#1}}


%hệ và tuyển phương trình  
\newcommand{\hhpt}[2]{\left\lbrace\begin{array}{l}#1\\ #2\end{array}\right.}
\newcommand{\hbpt}[3]{\left\lbrace\begin{array}{l}#1\\ #2\\ #3 \end{array}\right.}
\newcommand{\hppt}[4]{\left\lbrace\begin{array}{l}#1\\ #2\\ #3\\ #4 \end{array}\right.}

\newcommand{\thpt}[2]{\left[ \begin{array}{l}#1\\#2\end{array}\right.}
\newcommand{\tbpt}[3]{\left[ \begin{array}{l}#1\#2\\ #3 \end{array}\right.}
\newcommand{\tppt}[4]{\left[ \begin{array}{l}#1\\ #2\\ #3\\ #4 \end{array}\right.}


\newcommand{\motdong}[4]{\begin{tabbing}\hspace{.05\textwidth }\=\hspace{.25\textwidth }\=\hspace{.25\textwidth }\=\hspace{.25\textwidth }\=\kill \> {\textbf{A.}} $#1$. \>  {\textbf{B.}} $#2$. \>  {\textbf{C.}} $#3$. \> {\textbf{D.}} $#4$.\end{tabbing}\vspace*{-.5cm} }

\newcommand{\haidong}[4]{

\begin{tabbing}\hspace{.05\textwidth }\=\hspace{.475\textwidth }\=\kill\> {\textbf{A.}} $#1$. \>  {\textbf{B.}} $#2$. \\ 
\> {\textbf{C.}} $#3$. \> {\textbf{D.}} $#4$.\end{tabbing} \vspace*{-.5cm} }




%bảng biến thiên
%\usepackage{luamplib}

\newcommand{\bbtcd}[2]{
\begin{mplibcode}
	input tableauVariation
	beginTableau(0) newLigneVariables(btex $x$ etex);val(btex $-\infty $ etex);
	val(btex $#1$ etex);
	val(btex $+\infty $ etex); 
	newLigneSignes(btex $ f'(x)$ etex);plus; val(btex 0 etex);moins;
	newLigneVariations(btex $f(x)$ etex);valPos(btex $-\infty$ etex,0); valPos(btex   $#2$
	etex,1);valPos(btex $-\infty$ etex,0);
	endTableau;
\end{mplibcode}
}


\newcommand{\bbtct}[2]{
\begin{mplibcode}
	input tableauVariation
	beginTableau(0) newLigneVariables(btex $x$ etex);val(btex $-\infty $ etex);
	val(btex $#1$ etex);
	val(btex $+\infty $ etex); 
	newLigneSignes(btex $ f'(x)$ etex);moins; val(btex 0 etex);plus;
	newLigneVariations(btex $f(x)$ etex);valPos(btex $+\infty$ etex,1); valPos(btex   $#2$
	etex,0);valPos(btex $+\infty$ etex,1);
	endTableau;
\end{mplibcode}
}

\newcommand{\bbtcdct}[4]{
\begin{mplibcode}
	input tableauVariation
	beginTableau(0) newLigneVariables(btex $x$ etex);val(btex $-\infty $ etex);
	val(btex $#1$ etex);	val(btex $#2$ etex);
	val(btex $+\infty $ etex); 
	newLigneSignes(btex $ f'(x)$ etex);plus; val(btex 0 etex);moins;
val(btex 0 etex);plus;
	newLigneVariations(btex $f(x)$ etex);valPos(btex $-\infty$ etex,0); valPos(btex   $#3$etex,1);valPos(btex $#4$ etex,0); 
	valPos(btex $+\infty$ etex,1);
	endTableau;
\end{mplibcode}
}



\newcommand{\bbtctcd}[4]{
\begin{mplibcode}
	input tableauVariation
	beginTableau(0) newLigneVariables(btex $x$ etex);val(btex $-\infty $ etex);
	val(btex $#1$ etex);	val(btex $#2$ etex);
	val(btex $+\infty $ etex); 
	newLigneSignes(btex $ f'(x)$ etex);moins; val(btex 0 etex);plus;
val(btex 0 etex);moins;
	newLigneVariations(btex $f(x)$ etex);valPos(btex $+\infty$ etex,1); valPos(btex   $#3$etex,0);valPos(btex $#4$ etex,1); 
	valPos(btex $-\infty$ etex,0);
	endTableau;
\end{mplibcode}
}



\newcommand{\bbtcdctcd}[6]{
\begin{mplibcode}
	input tableauVariation
	beginTableau(0) newLigneVariables(btex $x$ etex);val(btex $-\infty $ etex);
	val(btex $#1$ etex);	val(btex $#2$ etex);val(btex $#3$ etex);
	val(btex $+\infty $ etex); 
	newLigneSignes(btex $ f'(x)$ etex);plus; val(btex 0 etex);moins;
val(btex 0 etex);plus;val(btex 0 etex);moins;
	newLigneVariations(btex $f(x)$ etex);valPos(btex $-\infty$ etex,0); valPos(btex   $#4$etex,1);valPos(btex $#5$ etex,0); valPos(btex $#6$ etex,1); 
	valPos(btex $-\infty$ etex,0);
	endTableau;
\end{mplibcode}
}

\newcommand{\bbtctcdct}[6]{
\begin{mplibcode}
	input tableauVariation
	beginTableau(0) newLigneVariables(btex $x$ etex);val(btex $-\infty $ etex);
	val(btex $#1$ etex);	val(btex $#2$ etex);val(btex $#3$ etex);
	val(btex $+\infty $ etex); 
	newLigneSignes(btex $ f'(x)$ etex);moins; val(btex 0 etex);plus;
val(btex 0 etex);moins;val(btex 0 etex);plus;
	newLigneVariations(btex $f(x)$ etex);valPos(btex $+\infty$ etex,1); valPos(btex   $#4$etex,0);valPos(btex $#5$ etex,1); valPos(btex $#6$ etex,0); 
	valPos(btex $+\infty$ etex,1);
	endTableau;
\end{mplibcode}
}


\newcommand{\bbtnbt}[2]{
\begin{mplibcode}
input tableauVariation
beginTableau(0) newLigneVariables(btex $x$ etex);
val(btex $-\infty $ etex);val(btex $#1$ etex);val(btex $+\infty $ etex);
 newLigneSignes(btex $\atop{\displaystyle f'(x)}$ etex);plus; nonDefBarre;plus;
newLigneVariations(btex $\atop{\displaystyle f(x)}$ etex);
valPos(btex $#2$ etex,0); limGauche(btex $+\infty $ etex,1);nonDefBarre;limDroite(btex $-\infty$ etex,0);valPos(btex $#2$  etex,1);
endTableau;
\end{mplibcode}
}




\newcommand{\bbtnbg}[2]{
\begin{mplibcode}
input tableauVariation
beginTableau(0) newLigneVariables(btex $x$ etex);
val(btex $-\infty $ etex);val(btex $#1$ etex);val(btex $+\infty $ etex);
 newLigneSignes(btex $\atop{\displaystyle f'(x)}$ etex);moins; nonDefBarre;moins;
newLigneVariations(btex $\atop{\displaystyle f(x)}$ etex);
valPos(btex $#2$ etex,1); limGauche(btex $-\infty $ etex,0);nonDefBarre;limDroite(btex $+\infty$ etex,1);valPos(btex $#2$  etex,0);
endTableau;
\end{mplibcode}
}

\newcommand{\bbtcdkxdct}[5]{
\begin{mplibcode}
input tableauVariation
beginTableau(0) newLigneVariables(btex $x$ etex);
val(btex $-\infty $ etex);val(btex $#1$ etex);val(btex $#2$ etex);val(btex $#3$ etex);val(btex $+\infty $ etex);
 newLigneSignes(btex $f'(x)$ etex);plus; val(btex 0 etex);moins;nonDefBarre;moins;val(btex 0 etex);plus;
newLigneVariations(btex $f(x)$ etex);
valPos(btex $-\infty$ etex,0); valPos(btex $#4$  etex,1);limGauche(btex $-\infty $ etex,0);nonDefBarre;limDroite(btex $+\infty$ etex,1);valPos(btex $#5$  etex,0);valPos(btex $+\infty$  etex,1);
endTableau;
\end{mplibcode}
}


\newcommand{\bbtctkxdcd}[5]{
\begin{mplibcode}
input tableauVariation
beginTableau(0) newLigneVariables(btex $x$ etex);
val(btex $-\infty $ etex);val(btex $#1$ etex);val(btex $#2$ etex);val(btex $#3$ etex);val(btex $+\infty $ etex);
 newLigneSignes(btex $f'(x)$ etex);moins; val(btex 0 etex);plus;nonDefBarre;plus;val(btex 0 etex);moins;
newLigneVariations(btex $f(x)$ etex);
valPos(btex $+\infty$ etex,1); valPos(btex $#4$  etex,0);limGauche(btex $+\infty $ etex,1);nonDefBarre;limDroite(btex $-\infty$ etex,0);valPos(btex $#5$  etex,1);valPos(btex $-\infty$  etex,0);
endTableau;
\end{mplibcode}
}

\newcommand{\bbtgachcheo}[4]{
\begin{mplibcode}
input tableauVariation
beginTableau(0)

newLigneVariables(btex $x$ etex);
val(btex $-\infty $ etex);
val(btex $#1$ etex);
val(btex $#2$ etex);
val(btex $+\infty $ etex);
 
newLigneSignes(btex $f'(x)$ etex);moins; debutNonDef;finNonDef;plus;
newLigneVariations(btex $ f(x)$ etex);valPos(btex $+\infty $ etex,1); limGauche(btex $#3$ etex,0);debutNonDef;finNonDef;limDroite(btex $#4$ etex,0);valPos(btex $+\infty $ etex,1);
endTableau;
\end{mplibcode}
}

\newcounter{vidu}\setcounter{vidu}{0}
\newcounter{baitap}\setcounter{baitap}{0}
\newcommand{\baitap}{\refstepcounter{baitap}Bài tập \thebaitap.\quad }
\newcommand{\vidu}{\refstepcounter{vidu}Ví dụ \thevidu.\quad }

\def\dis{\displaystyle}

\newcommand{\matphang}[1]{
\begin{mplibcode}
input couleur;
beginfig(1);
numeric u;
pair t,r;
transform T,S;
path p;

u= 1cm;
t=(4u,0u); r=(1u,2u);
T = identity shifted t;
S = identity shifted r;


z0=(0u,0u);
z1 = z0 transformed T;
z2 = z0 transformed S;
z3 = z0 transformed T transformed S;


p:=z0--z1--z3--z2--cycle;

fill p withcolor bleu_ciel;


draw z0--z2;
draw z2--z3;


pickup pencircle scaled 2pt;
draw z0--z1;
draw z1--z3;


pickup pencircle scaled 0.5pt;


label.urt(btex $#1$ etex, z0+(0.1u,0u));
endfig;
end
\end{mplibcode}
}


\newcommand{\giaotuyen}[3]{
\begin{mplibcode}
input couleur;
beginfig(2);

numeric u;
pair t,r;
transform T,S;
path p[],q[];

u= 1cm;


t=(4u,0u); r=(1u,2u);
T = identity shifted t;
S = identity shifted r;

z0=(0u,0u);
z1 = z0 transformed T;
z2 = z0 transformed S;
z3 = z0 transformed T transformed S;

path p;
p:=z0--z1--z3--z2--cycle;
fill p withcolor bleu_ciel;



draw z0--z2;
pickup pencircle scaled 2pt;
draw z0--z1;
draw z1--z3;
pickup pencircle scaled 0.5pt;
label.urt(btex $#1$ etex, z0+(0.1u,0u));


z5 = 0.5[z0,z1];
z6 = z5 transformed S;
z7 = (2.5u,-1u);
z8 = z7 transformed S;
z5 = 0.5[z9,z7];
z10= z9 transformed S;
p1 = z2--z6;
q1= z9--z10;
z11 = p1 intersectionpoint q1;
p2 = z7--z8;
q2= z5--z1;
z12 = p2 intersectionpoint q2;


path pp;
pp=z7--z8--z10--z9--cycle;
fill pp withcolor green;



path ppp;
ppp:=z5--z6--z11--z9--cycle;
fill ppp withcolor (green+.55*bleu_ciel);


path Pp;
Pp:=z5--z12--z8--z6--cycle;
fill Pp withcolor (green +.75*bleu_ciel);

draw z5--z6 withcolor bleu; 
draw z9--z7;
draw z9--z10;
draw z10--z6;
draw z6--z3;
draw z6--z8 dashed evenly;

draw z11--z6 dashed evenly;
draw z2--z11;
draw z5--z12 withpen pencircle scaled 2pt;
draw z12--z8 dashed evenly;
draw z7--z12;
label.rt(btex $#2$ etex, z9);


label.rt(btex $#3$ etex, 0.5[z5,z6])withcolor bleu;
endfig;
end
\end{mplibcode}
}



\newcommand{\mpABC}[4]{
\begin{mplibcode}
input couleur;
beginfig(3);
numeric u;
pair t,r;
transform T,S;

u= 1cm;
t=(4u,0u); r=(1u,2u);
T = identity shifted t;
S = identity shifted r;


z0=(0u,0u);
z1 = z0 transformed T;
z2 = z0 transformed S;
z3 = z0 transformed T transformed S;


path p;
p:=z0--z1--z3--z2--cycle;
fill p withcolor bleu_ciel;


draw z0--z2;
draw z2--z3;
pickup pencircle scaled 2pt;
draw z0--z1;
draw z1--z3;
pickup pencircle scaled 0.5pt;
label.urt(btex $#1$ etex, z0+(0.1u,0u));


dotlabel.top(btex $#2$ etex, (2u,1.5u));
dotlabel.top(btex $#3$ etex, (1u,.5u));
dotlabel.top(btex $#4$ etex, (3u,1u));
endfig;
end
\end{mplibcode}
}



\newcommand{\mpPddcn}[3]{
\begin{mplibcode}
input couleur;
beginfig(4);
numeric u;
pair t,r;
transform T,S;

u= 1cm;
t=(4u,0u); r=(1u,2u);
T = identity shifted t;
S = identity shifted r;


z0=(0u,0u);
z1 = z0 transformed T;
z2 = z0 transformed S;
z3 = z0 transformed T transformed S;




path p;
p:=z0--z1--z3--z2--cycle;
fill p withcolor bleu_ciel;




draw z0--z2;
draw z2--z3;
pickup pencircle scaled 2pt;
draw z0--z1;
draw z1--z3;
pickup pencircle scaled 0.5pt;
label.urt(btex $#1$ etex, z0+(0.1u,0u));


z4 = (0.8u,1u); z5 =(4u,1.5u);draw z4--z5;
z6 = (1u,1.6u);z7=(3.6u,0.4u);draw z6--z7;
label.rt(btex $#2$ etex, z5);
label.rt(btex $#3$ etex, z7);
endfig;
end
\end{mplibcode}
}

\newcommand{\mpPAd}[3]{
\begin{mplibcode}
input couleur;
beginfig(5);

numeric u;
pair t,r;
transform T,S;

u= 1cm;
t=(4u,0u); r=(1u,2u);
T = identity shifted t;
S = identity shifted r;


z0=(0u,0u);
z1 = z0 transformed T;
z2 = z0 transformed S;
z3 = z0 transformed T transformed S;


path p;
p:=z0--z1--z3--z2--cycle;
fill p withcolor bleu_ciel;


draw z0--z2;
draw z2--z3;
pickup pencircle scaled 2pt;
draw z0--z1;
draw z1--z3;
pickup pencircle scaled 0.5pt;
label.urt(btex $#1$ etex, z0+(0.1u,0u));


dotlabel.top(btex $#2$ etex, (2u,1.5u));
z4 = (1u,.5u);
z5 = (3u,1u);
draw z4--z5;
label.top(btex $#3$ etex, z5);
endfig;

end
\end{mplibcode}
}


\newcommand{\mpPddss}[3]{
\begin{mplibcode}
input couleur;
beginfig(6);

numeric u;
pair t,r;
transform T,S;

u= 1cm;
t=(4u,0u); r=(1u,2u);
T = identity shifted t;
S = identity shifted r;


z0=(0u,0u);
z1 = z0 transformed T;
z2 = z0 transformed S;
z3 = z0 transformed T transformed S;



path p;
p:=z0--z1--z3--z2--cycle;
fill p withcolor bleu_ciel;



draw z0--z2;
draw z2--z3;
pickup pencircle scaled 2pt;
draw z0--z1;
draw z1--z3;
pickup pencircle scaled 0.5pt;
label.urt(btex $#1$ etex, z0+(0.1u,0u));


z4 = (1u,1.3u); z5 =(3u,1.8u);draw z4--z5;
z6 = (1.5u,0.2u);z7 = z6 shifted (z5-z4);draw z6--z7;
label.rt(btex $#2$ etex, z5);
label.rt(btex $#3$ etex, z7);
endfig;
end
\end{mplibcode}
}




\newcommand{\cheonhauPddA}[4]{
\begin{mplibcode}
input couleur;
beginfig(7);
%representation de deux droites non parallelles et sans point commun
numeric u;
pair t,r;
transform T,S;

u= 1cm;
t=(4u,0u); r=(1u,2u);
T = identity shifted t;
S = identity shifted r;
path p[];

%tracé du plan
z0=(0u,0u);
z1 = z0 transformed T;
z2 = z0 transformed S;
z3 = z0 transformed T transformed S;


path p;
p:=z0--z1--z3--z2--cycle;
fill p withcolor bleu_ciel;


draw z0--z2;
draw z2--z3;
pickup pencircle scaled 2pt;
draw z0--z1;
draw z1--z3;
pickup pencircle scaled 0.5pt;
label.urt(btex $#1$ etex, z0+(0.1u,0u));

%tracé de la droite
z4 = (0.5u,.5u);
z5 = (4u,1u);
draw z4--z5;
label.top(btex $#2$ etex, z5);

%tracé du point
z6 = (2u,1.5u);
z7 = (0.2u,1u);
dotlabel.rt(btex $#4$ etex, z6);

%on met en place les pointillés
z9 = z6 shifted 2z7;
z10 = z6 shifted -3z7;
p1 = z6--z10;
p2 = z0--z1;
z11 = p1 intersectionpoint p2;
draw z6--z11 dashed evenly;
draw z11--z10;
draw z6--z9;

%le nom de la droite
label.rt(btex $#3$ etex, z10);
endfig;
end
\end{mplibcode}
}


\newcommand{\gtmpssPQdd}[4]{
\begin{mplibcode}
input couleur;
beginfig(8);

numeric u;
pair t,r;
transform T,S;
path p[],q[];
u= 1cm;
t=(4u,0u); r=(1u,2u);
T = identity shifted t;
S = identity shifted r;


z0=(0u,0u);
z1 = z0 transformed T;
z2 = z0 transformed S;
z3 = z0 transformed T transformed S;



path p;
p:=z0--z1--z3--z2--cycle;
fill p withcolor bleu_ciel;



draw z0--z2;
pickup pencircle scaled 2pt;
draw z0--z1;
draw z1--z3;
pickup pencircle scaled 0.5pt;
label.urt(btex $#1$ etex, z0+(0.1u,0u));


z4=(0u,-2.2u);
z5 = z4 transformed T;
z6 = z4 transformed S;
z7 = z4 transformed T transformed S;


path pp;
pp:=z4--z5--z7--z6--cycle;
fill pp withcolor 1.2*green;



draw z4--z6;
pickup pencircle scaled 2pt;
draw z4--z5;
draw z5--z7;
pickup pencircle scaled 0.5pt;
label.urt(btex $#2$ etex, z4+(0.1u,0u));


z8=1/5[z0,z1]+1/2[z0,z2];
z9=3/5[z4,z5]+1/2[z6,z4];

z10 = z8 transformed S;

z11 = z9 transformed S;
p1 = z10--z11;
q1 = z2--z3;
q2 = z0--z1;
q3 = z6--z7;
z12 = p1 intersectionpoint q1;
z13 = p1 intersectionpoint q2;
z14 = p1 intersectionpoint q3;

p2 = z8--z10;
z15 = p2 intersectionpoint q1;

p3 = z8--z9;
z16 = p3 intersectionpoint q2;
z17 = p3 intersectionpoint q3;
q4 = z4--z5;
z18 = p3 intersectionpoint q4;

p4 = z9--z11;
z19 = p4 intersectionpoint q4;


path k[];

k[1]:=z8--z10--z11--z9--cycle;
fill k[1] withcolor rose;


k[2]:=z8--z15--z12--z16--cycle;
fill k[2] withcolor (.75rose+.15bleu_ciel);



k[3]:=z16--z13--z12--cycle;
fill k[3] withcolor (.75rose+.35bleu_ciel);


k[4]:=z18--z17--z14--z19--cycle;
fill k[4] withcolor (.75rose+.25green);


k[5]:=z18--z19--z11--z14--cycle;
fill k[5] withcolor (.45rose+.45green);


draw z9--z19;draw z19--z11 dashed evenly;


draw z8--z9;
draw z8--z10;
draw z10--z12; draw z12--z13 dashed evenly;draw z13--z14;draw z14--z11 dashed evenly;
draw z2--z15; draw z15--z12 dashed evenly;draw z12--z3;
draw z16--z12;
draw z6--z17;draw z17--z14 dashed evenly;draw z14--z7;
draw z18--z14;

pickup pencircle scaled 2pt;
draw z4--z5;
draw z0--z1;


label.rt(btex $#3$ etex, 0.5[z12,z16]);
label.rt(btex $#4$ etex, 0.5[z18,z14]);
endfig;

end
\end{mplibcode}
}



\newcommand{\haimpssdtPQdd}[4]{
\begin{mplibcode}
input couleur;
beginfig(9);

numeric u;
pair t,r;
transform T,S;
path p[];

u= 1cm;
t=(4u,1u); r=(0u,2u);
T = identity shifted t;
S = identity shifted r;


z0=(0u,0u);
z1 = z0 transformed T;
z2 = z0 transformed S;
z3 = z0 transformed T transformed S;


path p;
p:=z0--z1--z3--z2--cycle;
fill p withcolor bleu_ciel;


draw z0--z2;
draw z2--z3;
draw z1--z3;
pickup pencircle scaled 0.5pt;
label.urt(btex $#1$ etex, z0+(0.1u,0u));


z4=0.5[z0,z1];
z5 = z4 transformed S;

z6 = z4 shifted (1.5u,-1u);
z4=0.5[z6,z7];
z8 = z6 transformed S;
z9 = z7 transformed S;

p1= z4--z1;
p2 = z6--z8;
z10 = p1 intersectionpoint p2;

p3 = z2--z3; p4 = z9--z7;
z11 = p3 intersectionpoint p4;



path k[];

k[1]:=z8--z9--z7--z6--cycle;
fill k[1] withcolor green;


k[2]:=z11--z5--z4--z7--cycle;
fill k[2] withcolor (.75bleu_ciel+green);



k[3]:=z5--z8--z10--z4--cycle;
fill k[3] withcolor (green+.55bleu_ciel);

draw z5--z11;


draw z11--z9;draw z11--z7 dashed evenly; draw z7--z4 dashed evenly;
label.lft(btex $#2$ etex, z6+(0u,0.5u));
draw z4--z5 withcolor red;
label.lft(btex $#3$ etex,0.5[z4,z5])withcolor red;
draw z0--z4;draw z4--z6;draw z6--z8;draw z8--z9;draw z4--z7 dashed evenly;
draw z4--z10 dashed evenly;draw z10--z1;


z12 = (-1u,0u);
z13 = z12 shifted (0u,3u);
draw z13--z12 withcolor red;
label.lft(btex $#4$ etex, z12)withcolor red;
endfig;
end
\end{mplibcode}
}


\newcommand{\haimpdtssddd}[3]{
\begin{mplibcode}
input couleur;
beginfig(10);

numeric u;
pair t,r,v;
transform T,S,V;
path p[];

u= 1cm;
t=(3u,2u); r=(-2u,1.5u);v=(-3u,0u); 
T = identity shifted t;
S = identity shifted r;
V = identity shifted v;
z0=(0u,0u);
z1 = z0 transformed T;
z2 = z0 transformed S;
z3 = z0 transformed S transformed T;
z4 = z0 transformed V;
z5 = z0 transformed V transformed T;

z101=0.1[z0,z1]; z10=0.9[z0,z1];
z23 = z101 transformed S; z32=z10 transformed S;
z45 = z101 transformed V; z54=z10 transformed V;

p1 = z101--z23; p2 = z4--z5;
z11 = p1 intersectionpoint p2;


path k[];

k[1]:=z45--z54--z32--z23--cycle;
fill k[1] withcolor green;



k[3]:=z23--z101--z10--z32--cycle;
fill k[3] withcolor bleu_ciel;

k[2]:=z23--z32--z54--z11--cycle;
fill k[2] withcolor (.75bleu_ciel+green);







draw z0--z1 ;
draw z2--z3 withcolor red;



draw z101--z23 ; draw z10--z32 ;
draw z23--z45 ; draw z32--z54 dashed evenly ;

draw z4--z11 ;
draw z11--z5 dashed evenly ;

label.rt(btex $#1$ etex, z0);
label.rt(btex $#2$ etex, z4);
label.top(btex $#3$ etex, z2);

endfig;
end
\end{mplibcode}
}




\newcommand{\haimpssPddQdd}[6]{
\begin{mplibcode}
input couleur;
beginfig(11);

numeric u;
pair t,r;
transform T,S;

u= 1cm;
t=(4u,0u); r=(1u,2u);
T = identity shifted t;
S = identity shifted r;

%premier plan P
z0=(0u,0u);
z1 = z0 transformed T;
z2 = z0 transformed S;
z3 = z0 transformed T transformed S;


path p;
p:=z0--z1--z3--z2--cycle;
fill p withcolor bleu_ciel;


draw z0--z2;
draw z2--z3;
pickup pencircle scaled 2pt;
draw z0--z1;
draw z1--z3;
pickup pencircle scaled 0.5pt;
label.urt(btex $#1$ etex, z0+(0.1u,0u));


z4 = (0.8u,1u); z5 =(4u,1.5u);draw z4--z5;
z6 = (1u,1.6u);z7=(3.6u,0.4u);draw z6--z7;
label.rt(btex $#2$ etex, z5);
label.rt(btex $#3$ etex, z7);


z10=(0u,-2.5u);
z11 = z10 transformed T;
z12 = z10 transformed S;
z13 = z10 transformed T transformed S;


path pP;
pP:=z10--z11--z13--z12--cycle;
fill pP withcolor green;


draw z10--z12;
draw z12--z13;
pickup pencircle scaled 2pt;
draw z10--z11;
draw z11--z13;
pickup pencircle scaled 0.5pt;
label.urt(btex $#4$ etex, z10+(0.1u,0u));


z14 = (0.8u,-1.5u); z15 =(4u,-1u);draw z14--z15;
z16 = (1u,-0.9u);z17=(3.6u,-2.1u);draw z16--z17;
label.rt(btex $#5$ etex, z15);
label.rt(btex $#6$ etex, z17);
endfig;
end
\end{mplibcode}
}

\newcommand{\dtssmpdPd}[3]{
\begin{mplibcode}
input couleur;
beginfig(12);

numeric u;
pair t,r;
transform T,S;

u= 1cm;
t=(4u,0u); r=(1u,2u);
T = identity shifted t;
S = identity shifted r;


z0=(0u,0u);
z1 = z0 transformed T;
z2 = z0 transformed S;
z3 = z0 transformed T transformed S;


path p;
p:=z0--z1--z3--z2--cycle;
fill p withcolor bleu_ciel;


draw z0--z2;
draw z2--z3;
pickup pencircle scaled 2pt;
draw z0--z1;
draw z1--z3;
pickup pencircle scaled 0.5pt;
label.urt(btex $#2$ etex, z0+(0.1u,0u));


z4 = (1u,.5u);
z5 = (4u,1u);
draw z4--z5;
label.top(btex $#3$ etex, z4);

%la droite en dehors du plan
z6 = (1u,2.5u);
z7 = z6 shifted z5-z4;
draw z6--z7;
label.top(btex $#1$ etex, z6);
endfig;
end
\end{mplibcode}
}


\def\today{Ngày \number\day\ tháng \number\month\ năm \number\year\ }